\section{导读：这不是人物稿，而是一条研究路线}

这期访谈长达六小时四十五分钟，表面上是谢赛宁第一次系统讲述自己的成长、求学、研究、创业和价值判断；更深层的线索，是一个做 computer vision 和 representation learning 的研究者，如何在 LLM 叙事压倒一切的时刻，重新定义“视觉”“世界模型”和“机器人大脑”。因此，这份笔记不会按流水账逐字复述，而是把访谈整理成一条可学习的研究路线。

本期有三条主线。第一条是个人路线：从上海交大 ACM 班、UCSD、FAIR、DeepMind、NYU 到 AMI Labs，每次选择都围绕“和谁做、做什么问题、是否有足够资源”。第二条是技术路线：从视觉任务、层次化表征、视频理解到 predictive world model。第三条是组织路线：当研究问题需要远超论文实验室的资源，又不适合被产品军备竞赛完全吞没时，创业公司要怎样保留 research oxygen。

\begin{importantbox}{本期核心命题}
谢赛宁在访谈中反复强调：LLM 是智能系统的重要组成部分，但不是全部。真正面向物理世界的智能，需要从连续、高维、有噪声的信号中学习任务相关表征，预测 action 的后果，并把预测用于 planning、decision making 和安全控制。
\end{importantbox}

\begin{knowledgebox}{视觉策略说明}
本视频是固定访谈画面，没有教学 slides、白板、产品演示或可读图表。按本仓库播客标准，正文不重复插入人物帧；封面用于来源识别，正文用世界模型闭环、能力阶梯、组织模型和概念表来承载教学内容。
\end{knowledgebox}

\subsection{本章小结}

这期节目的价值，不在于知道谢赛宁拒绝过谁、加入过哪里，而在于理解一个技术观点如何经过个人经历、研究品味和组织条件逐渐成形。后文会把访谈拆成：研究轨迹、世界模型定义、LLM 与物理世界的边界、AMI 的组织命题，以及结尾关于智能和哲学引用的反思。

